\scanpage{13}
la deuxième \deleted{\unsure{se}} se ramène à:
\[
E_n={}_{f^n}E_{n+1}
\]
qui par $a_n)$ équivaut à $\supset$, donc à ceci

soit $x\in E_n$, $h\in A$, \deleted{$\unsure{f^nx}$} tels qu'il existe $\lambda\in A$, avec
\[
f^n(x,h)=\lambda(\xi_n,f)
\]
alors on a $h\in fA$. En effet, on aura
\[
f^nh=\lambda f\quad\text{d'où}\quad\lambda=f^{n-1}h,
\qquad(\text{\uline{car $f$ est $A$-régulier}})
\]
puis $f^nx=\lambda\xi_n$ \deleted{\unsure{ou $f^n$}} i.e.
\[
0=f^{n-1}h\xi_n,\quad\text{or}\quad
(\overbrace{f^{n-1}\xi_n}^{\eta_n})\in{}_fE_n\quad\text{par }a_n),
\]
\deleted{d'où} et ${}_fE_n\simeq A/fA=B$ par $b_n)$, \deleted{donc}
enfin par construction on a $\xi_n\notin{}_{f^{n-1}}E_n$ i.e.
\[
\deleted{\unsure{f^n}}\eta_n\ne0,\quad\text{donc}\quad
h\eta_n\unsure{=|\ne}0\ \Rightarrow\ h\in fA,\quad\text{cqfd.}
\]

Enfin, vérifions $c_n$, on a
\[
\begin{aligned}
E_{n+1}/fE_{n+1}
&=(E_n\times A)/f(E_n\times A)+A(\xi_n,f)\\
&\simeq(E_n/fE_n\times A/fA)/A(\bar\xi_n,\bar f)\\
&\simeq([B\times(B/gB)^{n-1}]\times B)/B(g,0,0)\\
&\simeq B/gB\times(B/gB)^{n-1}\times B\simeq B\times(B/gB)^n.
\end{aligned}
\]
\margin{12}

